\chapter{Optimization}
\label{ch:optimization}

\noindent\textsc{Statistics is mostly optimization.} Find the parameter values that make the data most plausible. Find the line that minimizes the sum of squared errors. Find the model that maximizes the posterior probability. Once you understand optimization, you understand the engine that drives nearly every estimator in this book.

This chapter develops the calculus of finding maxima and minima of functions, building on the derivatives chapter and connecting directly to MLE, OLS, and beyond.

\section{Critical points}

For a function $f(x)$ defined on an interval, a \emph{critical point} is a value $x^*$ where either $f'(x^*) = 0$ or $f'$ does not exist. At a maximum or minimum on the interior of the interval, $f'$ must vanish (assuming $f$ is differentiable). This is \emph{Fermat's theorem}.

The intuition: at a peak, the function has just stopped going up and is about to start going down. The slope at the very top is zero. Similarly for a valley.

\paragraph{The first-order condition.} To find candidate maxima or minima:
\begin{equation}
  f'(x) = 0.
  \label{eq:foc-univariate}
\end{equation}
Solve this equation. Each solution is a critical point. Some will be maxima, some minima, some neither (saddle-like inflections).

\paragraph{The second-order condition.} To classify a critical point $x^*$:
\begin{itemize}[noitemsep]
\item If $f''(x^*) > 0$: local minimum (concave up)
\item If $f''(x^*) < 0$: local maximum (concave down)
\item If $f''(x^*) = 0$: test is inconclusive
\end{itemize}

\paragraph{Worked example.} Find the extrema of $f(x) = x^3 - 3x + 1$ on $[-2, 2]$.

First-order: $f'(x) = 3x^2 - 3 = 0 \Rightarrow x = \pm 1$.

Second-order: $f''(x) = 6x$. At $x = 1$, $f''(1) = 6 > 0$, so local minimum. At $x = -1$, $f''(-1) = -6 < 0$, so local maximum.



\begin{figure}[ht]
\centering
\includegraphics[width=\linewidth]{fig_optimization.pdf}
\caption{Left: the function $f(x) = x^3 - 3x + 1$ has critical points where $f'(x) = 0$. The second derivative classifies them: $f''(-1) = -6 < 0$ (max), $f''(1) = 6 > 0$ (min). Right: gradient descent on a quadratic. The path follows the negative gradient toward the minimum, converging in just a few steps.}
\label{fig:optimization}
\end{figure}

Function values: $f(-2) = -1$, $f(-1) = 3$ (local max), $f(1) = -1$ (local min), $f(2) = 3$. The global max on $[-2, 2]$ is 3 (achieved at both $x = -1$ and $x = 2$); the global min is $-1$ (at both $x = -2$ and $x = 1$).

\section{Multivariable optimization}

For a function of several variables, $f(x_1, \ldots, x_n)$, the analogue of $f'$ is the \emph{gradient}:
\[
  \nabla f = \left(\frac{\partial f}{\partial x_1}, \ldots, \frac{\partial f}{\partial x_n}\right).
\]
The first-order condition is $\nabla f = \bm{0}$: every partial derivative is zero. This produces a system of $n$ equations in $n$ unknowns.

The analogue of $f''$ is the \emph{Hessian matrix} of second partials:
\[
  \bm{H} = \begin{pmatrix}
    \frac{\partial^2 f}{\partial x_1^2} & \frac{\partial^2 f}{\partial x_1 \partial x_2} & \cdots \\
    \frac{\partial^2 f}{\partial x_2 \partial x_1} & \frac{\partial^2 f}{\partial x_2^2} & \cdots \\
    \vdots & \vdots & \ddots
  \end{pmatrix}.
\]
At a critical point, classifying maxima vs.\ minima requires examining the Hessian's eigenvalues:
\begin{itemize}[noitemsep]
\item All eigenvalues positive (Hessian is positive definite): local minimum.
\item All eigenvalues negative (negative definite): local maximum.
\item Mixed signs: saddle point (max in some directions, min in others).
\end{itemize}

\section{The OLS optimization, in full}

Now we can do the regression derivation properly. Minimize
\[
  S(\bbeta) = \|\bm{y} - \bm{X}\bbeta\|^2 = (\bm{y} - \bm{X}\bbeta)^\top(\bm{y} - \bm{X}\bbeta).
\]
Expand:
\[
  S(\bbeta) = \bm{y}^\top \bm{y} - 2\bbeta^\top \bm{X}^\top \bm{y} + \bbeta^\top \bm{X}^\top \bm{X} \bbeta.
\]
The gradient with respect to $\bbeta$:
\[
  \nabla S(\bbeta) = -2 \bm{X}^\top \bm{y} + 2 \bm{X}^\top \bm{X} \bbeta.
\]
Setting to zero:
\[
  \bm{X}^\top \bm{X} \hat\bbeta = \bm{X}^\top \bm{y}, \qquad \hat\bbeta = (\bm{X}^\top \bm{X})^{-1} \bm{X}^\top \bm{y}.
\]
The Hessian is $\bm{H} = 2\bm{X}^\top \bm{X}$, which is positive semi-definite (and positive definite if $\bm{X}$ has full column rank). So this critical point is a global minimum, and we have found the OLS estimator.

\section{The likelihood optimization (preview of MLE)}

Maximum likelihood estimation finds the parameter values that make the observed data most probable. Given a model $f(x; \theta)$ and data $x_1, \ldots, x_n$ assumed independent, the \emph{likelihood function} is
\begin{equation}
  L(\theta) = \prod_{i=1}^{n} f(x_i; \theta).
  \label{eq:likelihood}
\end{equation}
Maximizing this is mathematically equivalent to maximizing the \emph{log-likelihood}
\begin{equation}
  \ell(\theta) = \sum_{i=1}^{n} \log f(x_i; \theta).
  \label{eq:log-likelihood}
\end{equation}
The log turns a product into a sum, which is much easier to differentiate. The maximum likelihood estimator is
\[
  \hat\theta_{\text{MLE}} = \arg\max_\theta \ell(\theta).
\]
The first-order condition is the \emph{score equation}:
\[
  \frac{d\ell}{d\theta} = 0.
\]

\paragraph{Example: MLE for the normal mean.} Suppose $X_1, \ldots, X_n \sim N(\mu, \sigma^2)$ with $\sigma$ known. The log-likelihood is
\[
  \ell(\mu) = -\frac{n}{2}\log(2\pi\sigma^2) - \frac{1}{2\sigma^2}\sum_{i=1}^{n}(x_i - \mu)^2.
\]
Differentiate:
\[
  \frac{d\ell}{d\mu} = \frac{1}{\sigma^2}\sum_{i=1}^{n}(x_i - \mu) = 0.
\]
Solve: $\hat\mu = \bar{x}$.

The MLE for the normal mean is the sample mean. This isn't a coincidence; it's why the sample mean is the standard estimator.

\section{Constrained optimization (Lagrange multipliers)}

Sometimes you optimize subject to constraints. Minimize $f(x, y)$ subject to $g(x, y) = 0$. The Lagrangian is
\[
  \mathcal{L}(x, y, \lambda) = f(x, y) - \lambda g(x, y).
\]
First-order conditions:
\[
  \frac{\partial \mathcal{L}}{\partial x} = 0, \quad \frac{\partial \mathcal{L}}{\partial y} = 0, \quad \frac{\partial \mathcal{L}}{\partial \lambda} = 0.
\]
The third equation is just the constraint. The first two say that the gradient of $f$ is parallel to the gradient of $g$ at the optimum.

This is how economists derive consumer demand: maximize utility subject to a budget constraint. It's how regularized regression (ridge, lasso) is set up: minimize squared error subject to a constraint on the coefficient sizes.

\section{Numerical optimization}

For most real models, there's no closed-form solution. You compute the optimum numerically using iterative algorithms:

\paragraph{Gradient descent.} Start at $\theta_0$. Update:
\[
  \theta_{t+1} = \theta_t - \eta \cdot \nabla f(\theta_t).
\]
Take a step in the direction of steepest descent, scaled by a learning rate $\eta$.

\paragraph{Newton's method.} Use second-order information:
\[
  \theta_{t+1} = \theta_t - \bm{H}(\theta_t)^{-1} \nabla f(\theta_t).
\]
Faster convergence near the optimum, more expensive per step.

\paragraph{BFGS and L-BFGS.} Approximate the Hessian rather than computing it exactly. R's \texttt{optim()} default. Good for medium-sized problems.

The mathematics of why these converge (and when) is beautiful but beyond this chapter. For now: optimization is solved, you call a function, the function gives you the answer.

\begin{keyinsight}
Optimization is the engine that turns ``find the best parameter'' into ``solve this equation.''

OLS minimizes squared error. MLE maximizes log-likelihood. Posterior estimation maximizes the posterior. They differ in the objective function, not in the technique.

The first-order condition $\nabla f = 0$ is universal. The second-order condition tells you whether you found a max, a min, or a saddle.
\end{keyinsight}

\begin{codelab}
\begin{lstlisting}[language=Rlang]
# Find the maximum of f(x) = -x^2 + 4x - 3
f <- function(x) -x^2 + 4*x - 3
optimize(f, interval = c(-10, 10), maximum = TRUE)
# Should give x = 2, max = 1

# MLE for normal mean (sigma known)
set.seed(2024)
data <- rnorm(50, mean = 5, sd = 2)
nll <- function(mu) -sum(dnorm(data, mu, 2, log = TRUE))

# optim minimizes by default; we minimize negative log-likelihood
optim(par = 0, fn = nll, method = "Brent",
      lower = -10, upper = 10)
# Should give estimate close to 5 = the true mean
mean(data)  # for comparison
\end{lstlisting}
\end{codelab}

\section{Practice problems}

\begin{enumerate}
  \item Find the critical points of $f(x) = x^4 - 4x^3$ and classify each as max, min, or neither.
  \item For $f(x, y) = x^2 + xy + y^2 - 2x - 4y$, find the critical point by setting both partial derivatives to zero. Classify it using the Hessian.
  \item Show that the MLE of $p$ for $X_1, \ldots, X_n \sim \text{Bernoulli}(p)$ is the sample proportion $\bar{x}$. Write the log-likelihood, differentiate, and solve.
  \item Use Lagrange multipliers to find the rectangle of largest area with perimeter 100.
  \item In R, write a function that computes the negative log-likelihood for a normal sample, then use \texttt{optim} to find the joint MLE for $(\mu, \sigma)$.
  \item Show that for the OLS objective, the Hessian $2\bm{X}^\top \bm{X}$ is always positive semi-definite. (Hint: show that for any vector $\bm{v}$, $\bm{v}^\top \bm{X}^\top \bm{X} \bm{v} \geq 0$.)
\end{enumerate}

\begin{reflect}
$\triangleright$~State the first-order and second-order conditions for an optimum.

\smallskip
$\triangleright$~Why is OLS a special case of an optimization problem? What is being optimized, over what?

\smallskip
$\triangleright$~What's the relationship between MLE and OLS for a normal linear model?
\end{reflect}
